\documentclass[11pt]{article}
\usepackage[margin=1in]{geometry}
\usepackage{parskip}
\usepackage{hyperref}
\hypersetup{hidelinks}
\pagestyle{empty}

\begin{document}

\begin{flushright}
Panagiotis Karmiris\\
Independent Researcher, Greece\\
\href{mailto:unbinder@msn.com}{unbinder@msn.com}\\
ORCID: \href{https://orcid.org/0009-0007-7536-1467}{0009-0007-7536-1467}\\
14 September 2026
\end{flushright}

\vspace{1em}

Dear Editors of \textit{Physical Review D},

I am pleased to submit the manuscript entitled

\begin{quote}
\textbf{``A Matched Seven-Likelihood Test of a Zero-Particle-CDM Single-Metric Relativistic Modified-Gravity Cosmology''}
\end{quote}

for consideration in \textit{Physical Review D}.

The manuscript presents a matched cosmological likelihood test of a single-metric relativistic modified-gravity model in which the particle cold-dark-matter density and bare cosmological constant are both fixed to zero. The model combines an Aether Scalar--Tensor sector, a current-linked nonseparable scalar interaction, and a Type-II minimally modified gravity sector for late-time acceleration. It is tested against a matched flat-$\Lambda$CDM control using the same seven likelihood contributions from Planck 2018 high-$\ell$ TTTEEE-lite, low-$\ell$ TT, low-$\ell$ EE, Planck lensing, DESI DR2 BAO, Pantheon+, and SH0ES.

Both chains satisfy the common convergence criterion $R-1<0.01$. The lowest sampled independent seven-likelihood totals are $\chi^2_{\rm MG}=2441.26650$ and $\chi^2_{\Lambda\mathrm{CDM}}=2446.93250$, giving $\Delta\chi^2=-5.66600$. With seven fitted cosmological coordinates for modified gravity and six for the matched $\Lambda$CDM control, the corresponding differences are $\Delta\mathrm{AIC}=-3.66600$ and $\Delta\mathrm{BIC}=+2.11431$ under the stated bookkeeping convention. The improvement is concentrated in the Planck high-$\ell$, Pantheon+, and DESI DR2 sectors, while the SH0ES contribution itself is not responsible for the net gain.

The main significance of the work is therefore structural: a relativistic zero-particle-CDM realization remains statistically competitive with the matched $\Lambda$CDM control on this precision likelihood stack while using only one additional fitted cosmological coordinate. The manuscript also provides the action, background reduction, perturbative mapping, stability analysis, numerical implementation, likelihood construction, prior-boundary diagnostic, and reproducibility information needed to assess the result.

The development and reproducibility repository is available at
\href{https://github.com/KarmirisP/aest-nonseparable-cosmology}{github.com/KarmirisP/aest-nonseparable-cosmology}. The manuscript includes an explicit disclosure of the artificial-intelligence tools used for algebraic cross-checking, code generation and debugging, workflow organization, literature organization, and language editing; all reported numerical calculations were executed in the author's computing environment and the author assumes responsibility for the scientific content.

Thank you for considering this work for publication in \textit{Physical Review D}.

Sincerely,

\vspace{1em}
Panagiotis Karmiris

\end{document}
